\begin{figure}[ht]
\begin{AcademicBox}[案例 3：长上下文叙事理解（《基督山伯爵》）]
\small
\setlength{\parindent}{0pt}

% --- Context ---
\textbf{\textit{上下文：}} \\
《基督山伯爵》开篇长段落节选：
1810 年 2 月 24 日，马赛港外了望台发现“法老号”归航，船只经由多港返回并缓慢驶入港口；城内围观者因其异常航行状态而猜测是否发生不幸事件……
（原文篇幅较长，此处省略部分上下文）
\vspace{4pt} \hrule \vspace{4pt}

% --- Question ---
\textbf{\textit{问题：}} \\
是谁写信并将充满谎言的举报信送交检察官，从而陷害了埃德蒙·唐泰斯？
\vspace{4pt} \hrule \vspace{4pt}

% --- Correct Result (Full Width) ---
\textbf{\textcolor{teal}{正确预测（我们的方法 / 标准答案）：}} \hfill \textcolor{teal}{\checkmark} \\
\textbf{Danglars、Fernand 和 Caderousse。}
\vspace{4pt} \hrule \vspace{4pt}

% --- Incorrect Results (Grouped) ---
\textbf{\textcolor{red}{错误基线：}} \\
\footnotesize
\begin{itemize}[leftmargin=*, label={}, itemsep=4pt, topsep=2pt]
 \item \textbf{Llama3.1-8B-Instruct、DuoAttn、NSA：} \\
 \textbf{Fernand、Villefort 和 Danglars} \\
    \textit{（错误：实体混淆，把检察官 Villefort 误当成共谋者）}

 \item \textbf{InfLLM-V2：} \\
 \textbf{Fernand 和 Caderousse} \\
    \textit{（错误：信息不完整，遗漏主谋 Danglars）}

 \item \textbf{PruLong：} \\
 \textbf{Caderousse、Villefort 和 Fernand} \\
    \textit{（错误：实体混淆，错误包含 Villefort）}

\end{itemize}

\end{AcademicBox}
\caption{叙事实体追踪任务的定性比较。该任务要求识别陷害主角的具体共谋角色。我们的模型准确检索到正确三人组，而基线常将“Villefort（检察官）”幻觉性并入共谋者集合，未能区分策划者与后续司法角色。}
\label{fig:case_montecristo}
\end{figure}
